% !TeX root = ../wi-screen.tex

\chapter{引言}

经典逻辑非常有用，应用广泛，历史悠久，且相对简单。
但它有局限性：例如，它不能（且不可能）很好地处理自然语言中的某些表达方式，
如时态和虚拟语气，也无法处理某些结构，比如“Audrey 知道 $p$”。
它做出某些假设，例如，每个语句非真即假，且绝非既真又假。
它将某些公式判定为重言式、某些论证判定为有效，
尽管这些重言式和论证是对英语论证的形式化，
而有些人并不认为这些英语论证是真的或有效的——至少并非显然如此。
因此，似乎存在这样一些例子，表明经典逻辑的表达力不够充分，甚至在某些情况下出了错。

本书讨论一些替代性的、\emph{非经典}逻辑。
这些非经典逻辑要么比经典逻辑表达力更强，要么拥有不同的重言式或有效论证。
例如，时态逻辑通过表达时态的算子来扩展经典逻辑；
条件句逻辑拥有一个额外的、不同的条件句（“如果……那么”），
它不受所谓实质条件句悖论的困扰。
所有这些逻辑都通过新的算子或联结词来\emph{扩展}经典逻辑，
并归属于内涵逻辑这一广阔范畴。
其他逻辑，如多值逻辑、直觉主义逻辑和次协调逻辑，则拥有与经典逻辑相同的基本联结词，
但不同的推理被视为有效。
例如，在多值逻辑和直觉主义逻辑中，排中律 $!A \lor \lnot !A$ 不成立；
在次协调逻辑中，推理\emph{ex contradictione quodlibet}（矛盾推出任意命题），
即对任意~$!A$ 有 $\lfalse \Entails !A$，也不成立。

在\cref{part:classical}回顾经典命题逻辑基础之后，
我们将在\ref{part:many-valued}开始讨论非经典逻辑。
在那里，我们将放松经典逻辑所做的一个假设：一切非真即假。
我们有充分理由认为，英语中的某些语句既非真也非假——它们具有某种中间真值。
这样的例子包括涉及模糊性的语句（“Mary 很富有”），
其真假尚未确定的语句（“明天会有海战”），
以及——对哲学家来说很重要的——悖论性语句，如“这句话是假的”。
最早的非经典逻辑之一允许在经典的“真”和“假”之外还有其他真值；
它们被称为\emph{多值逻辑}。我们将在\cref{part:many-valued}讨论它们。

模态逻辑是通过算子 $\Box$（“方框”）和 $\Diamond$（“菱形”）对经典逻辑的扩展，
这些算子附着在公式上。
直观上，$\Box$ 可以读作“必然”，$\Diamond$ 读作“可能”，
因此 $\Box p$ 是“$p$ 必然为真”，而 $\Diamond p$ 是“$p$ 可能为真”。
由于必然性和可能性是根本的形而上学概念，模态逻辑显然具有重大的哲学意义。
它使得诸如“$\Box p \lif p$”（如果 $p$ 是必然的，那么它为真）
或“$\Diamond p \lif \Box\Diamond p$”（如果 $p$ 是可能的，那么它必然可能）
这样的形而上学原则得以形式化。

对于那种将 $\Box$ 解释为“必然”的逻辑，其语义相对简单：
不是给命题变元分配真值，而是解释~$\mModel{M}$ 给它们分配一个“世界”的集合
——直观上，这些是 $p$ 被解释为真的世界~$w$。
基于这样的解释，我们可以定义一个满足关系。
这个满足关系的定义使得 $\Box !A$ 在世界~$w$ 处被满足，当且仅当 $!A$ 在\emph{所有}世界处都被满足：$\mSat{M}{\Box !A}[w]$当且仅当$\mSat{M}{!A}[v]$对所有世界~$v$成立。
这对应着 Leibniz（莱布尼茨）的思想：必然为真的就是在所有可能世界中都为真的。

“必然”并非$\Box$算子的唯一解释方式，但它是标准解释——“必然”与“可能”即所谓的\emph{真势}模态。其他解释将$\Box$读作“（某个人~$A$）知道”、“某个人~$A$相信”、“应当如此”或“将永远为真”，这些分别是认知模态、信念模态、道义模态和时态模态。对$\Box$的不同解释会使不同的公式逻辑地为真，并判定不同的推理为有效。例如，所有必然为真的事物和所有已知为真的事物都是真的，因此$\Box !A \lif !A$在真势和认知解释下是逻辑真理。相反，并非所有被相信的事物或所有应当如此的事物实际上都如此，因此$\Box !A \lif !A$在信念或道义解释下不是逻辑真理。我们将在\cref{part:modal1,part:modal2}中一般性地讨论模态逻辑，在\cref{part:epistemic}中专门讨论认知逻辑。

为了处理模态算子的不同解释，语义通过一个世界间的关系——即所谓的可达关系——得到了扩展。此时，$\mSat{M}{\Box !A}[w]$，若$\mSat{M}{!A}[v]$对所有从~$w$可达的世界~$v$成立。由此产生的语义非常通用且强大，其基本思想可用于为基于其他内涵算子的逻辑提供语义解释。其中一种逻辑是模态逻辑的近亲，称为时态逻辑。它不止有一个模态算子$\Box$（加上它的对偶$\Diamond$），而是有诸如“总是$P$”、“$p$将为真”等\emph{时态算子}。我们将在\cref{part:temporal}中研究这些。

实质条件句最好解读为英语的直陈条件句（“如果$p$为真，那么$q$为真”），而虚拟条件句则使用虚拟语气：“如果$p$为真，那么$q$就会为真。”一个具有假前件的实质条件句是真的，但一个虚拟条件句则未必，例如，“如果人类有尾巴，他们就能飞。”在\cref{part:counterfactuals}中，我们讨论反事实条件句的逻辑。

直觉主义逻辑是一种基于 L. E. J. Brouwer（布劳威尔）的构造性数学分支的构造性逻辑。直觉主义逻辑因此具有哲学趣味——它在数学的构造性解释中扮演着重要角色——但也曾被二十世纪有影响力的英国哲学家 Michael Dummett（达米特）提出作为优于经典逻辑的逻辑。如上所述，直觉主义逻辑是非经典的，因为它有更少的有效推理和定理，例如，$!A \lor \lnot !A$ 和 $\lnot\lnot !A \lif !A$ 一般不成立。直观上，这是直觉主义原则的结果：除非你拥有某个东西的证明，否则它就不应算作真——你不应断言它。显然，存在我们既没有$!A$的证明也没有$\lnot !A$的证明的情况。直觉主义逻辑可以给出一种非常类似于模态逻辑的关系语义。我们将在\cref{part:intuitionistic}中讨论它。

经典逻辑最奇怪的特征之一是爆炸原理：从矛盾可以推出任何结论。这一原理源于我们设定经典逻辑语义的方式，但它\emph{非常}反直觉，并且违背了我们实际推理时的做法。毕竟，一旦你发现你相信的某些东西是矛盾的，你（通常）不会进而推断出任意主张，因为它们可以从你的信念中推出！这促使逻辑学家发展了一些逻辑系统，在这些系统中，\emph{ex contradictione, quodlibet}这一推理被阻断。其中一些仅仅是经典逻辑的进一步弱化，只是为了摆脱爆炸原理。另一些则更有哲学动机。爆炸原理之所以奇怪，部分原因是当它断定$!A$和$\lnot !A$一起后承$!B$时，前提和结论之间完全没有联系。难道在任何有效论证中不应该有这样的联系吗？前提难道不应该与结论\emph{相关}吗？这就引出了所谓的\emph{相干}（或\emph{相关}）逻辑。另一个动机是一种称为\emph{矛盾真理论}的哲学立场：相信可能存在真矛盾。（如果你相信矛盾可以为真，那么你大概不会希望它们后承任何东西，比如假命题。）所有这些逻辑都属于\emph{次协调逻辑}这一总称。我们将在\cref{part:paraconsistent}中讨论次协调逻辑。